\section{Related Work}

\subsection{O-RAN Architecture and Power-Outage Context}

O-RAN builds on the disaggregated 5G radio access network architecture in which centralized and distributed processing functions are separated from radio functions. In particular, 3GPP distinguishes Central Unit (CU) and Distributed Unit (DU) functions within the gNB architecture, while O-RAN specifies open and interoperable interfaces and functions around components such as the O-CU, O-DU, and O-RU~\cite{oranArchitecture,3gppTs38401}. The O-RU performs the radio-facing functions that enable wireless access, making its continued operation particularly important for users located within its service area.

In this analysis we consider a network-wide grid-power outage scenario in which the O-RUs operate using finite local battery capacity. The O-CU and O-DU are assumed to retain backup power and are therefore excluded from the modeled energy constraint. This abstraction, while not a complete representation of deployment-level power architecture, enables focused examination of the trade-off at the radio-access edge: maintaining an RU in an active state can sustain user connectivity, but it also depletes the stored energy required to maintain service as the outage persists.

\subsection{Sleep Modes and Energy-Aware Radio Access Management}

Sleep-mode operation has long been studied as a way to reduce the energy consumption of cellular base stations. Works such as ~\cite{wu2015sleep,lopezPerez2022ranEnergy} provide a broad survey of sleep-mode techniques, showing how radio resources can be selectively deactivated when full network capacity is unnecessary. Such techniques are commonly motivated by normal variations in traffic demand, where lightly loaded cells may enter a reduced-power state while nearby active cells absorb their traffic.

However, base-station sleep decisions influence coverage, user association, cell capacity, and service quality. Users served by a station entering sleep mode must be reassociated with an available active station~\cite{son2011energyDelay,oh2013dynamicSwitching}. Consequently, energy savings should be evaluated alongside their impact on network performance. During a power outage, the trade-off is more pronounced: conserving battery capacity is beneficial only if the remaining active RUs can sustain adequate emergency connectivity.

\subsection{Cellular Resilience During Power Outages}

Research on cellular resilience investigates how communication networks remain operational during disasters, infrastructure failures, and power outages. Existing strategies encompass distributed architectures, self-organizing control, backup-energy allocation, and adaptive recovery mechanisms~\cite{gomez2014disasterResilient,zhang2016selfOrganization,karaman2025resilientInfrastructure}. These studies establish that reliable emergency communication is contingent upon both radio-network functionality and the availability of energy resources.

Effective backup-battery management is critical when grid power is unavailable. Wang \emph{et al.}~\cite{wang2019backupBattery} demonstrate that battery capacity and allocation substantially influence service interruptions during cellular power outages. Unlike studies that optimize battery deployment, recovery mechanisms, or adaptive control, this analysis maintains a fixed battery configuration and concentrates on the operational impact of predefined RU sleep policies. While detailed simulators, emulators, and O-RAN experimentation platforms enable more realistic validation~\cite{patriciello2019e2e,nardini2020simu5g,bonatiOpenRANGym,bonati2021colosseum}, the simulator employed here offers a transparent, policy-level comparison within a standardized outage model while eliminating unnecessary complexities that do not benefit the simulation under the proposed model.

In summary, prior research demonstrates the benefits of sleep-mode operation and highlights the necessity of resilient communication during power outages. This study contributes by providing a comparative evaluation of three predefined RU policies: Always Active, Staggered Active, and Threshold Staggered Active. Instead of proposing an optimal controller, the analysis examines whether simple sleep schedules can simultaneously maintain emergency quality of service and extend RU battery lifetime under the specified outage conditions.
